\documentclass[11pt,a4paper]{article}
\usepackage[margin=23mm]{geometry}
\usepackage{amsmath,amssymb,booktabs,array}
\input{glyphtounicode}
\pdfgentounicode=1
\usepackage[hidelinks]{hyperref}
\hypersetup{hypertexnames=false}
\setlength{\parindent}{0pt}
\setlength{\parskip}{5pt}
\newcommand{\pos}[1]{[#1]_+}
\renewcommand{\labelitemi}{\(\bullet\)}
\begin{document}
\begin{center}
{\Large Physical qualification: what can be identified?}\\[4pt]
16 October 2026 --- MP6-Q1; no traffic experiment or parameter selection
\end{center}

\textbf{Status.} The selected conditional model is unchanged. Its forward
domain is qualified at WF-Q1. This note separates its physical meanings
and gives exact, conditional observation identities. It does not establish
an empirical parameter range, noisy-data identifiability or applied MP6.

\section*{1. A boundary recovery quantity is not a vehicle trajectory}
Retain $K=wk_j$, $c(s)=Ks/(w+s)$, $U(f)=wf/(K-f)$ and
$0<R\leq C(v_-)=c(v_-)<K$. In the selected node,
\begin{equation}
 c_E=1+\delta+\chi g_A\pos{u_A-u_E},\quad c_A=1,\quad
 B=R-p,\qquad g_A+c_Eg_E\leq B. \tag{1}
\end{equation}
The previously proved positive-budget domain makes clipping inactive.
Free marks are $u_j=v$; queued marks are $u_j=U(g_j)$. Their use as actual
entry/passage speed is an equilibrium approximation, not a consequence of
conservation. $\chi$ is fixed in physical units across advice.

For an isolated entry with mark $u\geq0$, define
\begin{align}
 S(\theta;u)&=\sqrt{(w+u)^2+2wa_r\theta},\quad
 V(\theta;u)=S(\theta;u)-w,\nonumber\\
 Q_{\rm iso}(\theta;u)&=K\left[\theta-
        \frac{S(\theta;u)-w-u}{a_r}\right],\nonumber\\
 \partial_\theta Q_{\rm iso}&=K(1-w/S)=c(V),\quad
 D(\theta;u)=\pos{R-c(V)}. \tag{2}
\end{align}
The derivative identity explains the retained profile. $V$ parametrizes
the boundary passing rate; it is not the admitted vehicle's speed.
For comparison, a vehicle accelerating constantly before a cap has speed
$u+a_r\theta$ along its own trajectory, which differs from $V$.
Actual acceleration observations constrain the source interpretation of
$a_r$; observed discharge tests (2). Fitting $V$ directly to measured vehicle
speeds would test a different assertion.

This distinction is explicit in the accompanying interpretation of equation
(1) of Leclercq et al. (2016). Their congested insertion-speed assumption in
Section 3.2 gives a restricted precedent for a flow-to-speed relation, not
a proof of every transient stream mark used here.

Units: $g_j,B,p,D,K$ are veh-eq/s; $u_j,w,V$ are m/s; $a_r$ is m/s$^2$;
$\delta$ is dimensionless; $\chi$ is s$^2$/(veh-eq m). The reference
mass $m=1$ veh-eq is physical normalization, not an integer service packet.

\newpage
\section*{2. What distinguishes recovery from current maneuver work?}
With the selected recent-mass renewal, $Y_E$ is cumulative admitted E and
$T(\xi)$ its mass-rank admission time:
\begin{equation}
 p(t)=\frac1m\int_{[Y_E(t)-m]_+}^{Y_E(t)}
 D(t-T(\xi);u_E(T(\xi)))\,d\xi. \tag{3}
\end{equation}
If E has finished entering and A remains queued, $g_E=0$ and its current
maneuver work vanishes. Work conservation in the selected node gives
\begin{equation}
 g_A=B=R-p,\qquad p=R-g_A. \tag{4}
\end{equation}
This identifies the account's value on that ideal saturated tail if $R$ is
independently known. It does not identify every profile parameter from one
tail, guarantee such a tail exists in the primary cohort, or make a
microscopic individual headway a continuous instantaneous rate. The later
execution proposal must specify observation scale and uncertainty.
If A is free, low $g_A$ alone does not identify $B$.

At a binding mixed-service observation with $g_E>0$, set
\begin{equation}
 X=g_A\pos{u_A-u_E},\qquad
 Z=\frac{R-p-g_A}{g_E}-1.
 \quad\Longrightarrow\quad Z=\delta+\chi X. \tag{5}
\end{equation}
After recovery and marks are constrained separately, two distinct covariate
values have determinant $X_2-X_1$ and give
\begin{equation}
 \chi=\frac{Z_2-Z_1}{X_2-X_1},\qquad
 \delta=\frac{X_2Z_1-X_1Z_2}{X_2-X_1}. \tag{6}
\end{equation}
Admissible estimates must still obey the selected nonnegative and
monotone-root conditions. Equal $X$ values do not separate the coefficients;
all $X=0$ gives no information about $\chi$. With nonbinding service,
the counterpart of (5) is only an inequality.

At a single observation without independent recovery information, the
accounting ambiguity is explicit:
\begin{equation}
 \delta^\prime=\delta+\varepsilon,\qquad
 p^\prime=p-\varepsilon g_E
 \quad\Longrightarrow\quad
 p^\prime+g_A+(1+\delta^\prime+\chi X)g_E
 =p+g_A+(1+\delta+\chi X)g_E. \tag{7}
\end{equation}
This is not a second admissible O-R trajectory or a global
non-identifiability theorem: the primed account need not satisfy (3).
It shows why discharge accounting alone does not establish the physical
partition. Distinct tails, mixtures and joint predictions are needed.

The profile already removes service through (3). Adding its integrated
loss again to $c_E$ would charge that same modeled loss twice.
Whether fitted immediate work represents a distinct real burden still
requires evidence; algebra alone cannot establish that physical fact.

\newpage
\section*{3. Which preparation and longitudinal scales are visible?}
For one constant advice, use traveled preparation distance $z$, remaining
ending density $e(z)$ and the selected H1 relation
\begin{equation}
 \frac{de}{dz}=-\frac1h\min\left\{e,
                  \frac{\Delta+e}{\gamma}\right\},\qquad
 \Delta=\frac{K}{w+v}-r. \tag{8}
\end{equation}
$h$ is an opportunity length and $\gamma$ is dimensionless receiving
penalty. On an informative donor-limited interval,
\begin{equation}
 e(z)=e(z_0)e^{-(z-z_0)/h},\qquad
 -\frac{d\log e}{dz}=\frac1h. \tag{9}
\end{equation}
On a receiving-limited interval with fixed $v,r$,
\begin{equation}
 e(z)+\Delta=[e(z_0)+\Delta]e^{-(z-z_0)/(\gamma h)},\qquad
 -\frac{d\log(e+\Delta)}{dz}=\frac1{\gamma h}. \tag{10}
\end{equation}
Thus that branch identifies a product. The transformation
$h^\prime=s h$, $\gamma^\prime=\gamma/s$ preserves (10), as long as the
branch and admissibility remain unchanged. Donor-limited observations can
constrain $h$ first. If no receiving-limited behavior is observed, the
penalty may only be bounded. A single final E count generally cannot
identify the whole law. These are ideal branch-conditional statements,
not statistical estimators for arbitrary noisy spatial observations.

For the earlier stationary correspondence, let $d$ be vehicle length plus
standing gap and $\tau_d$ a desired time gap. In the documented ideal
constant-speed Krauss case,
\begin{equation}
 k_{\rm ideal}(s)=\frac1{d+\tau_d s},\qquad
 f_{\rm ideal}(s)=\frac{s}{d+\tau_d s}. \tag{11}
\end{equation}
Choosing $k_j=1/d$ and $w=d/\tau_d$ gives
\begin{equation}
 \frac{K}{w+s}=\frac1{d+\tau_d s},\qquad
 c(s)=\frac{s}{d+\tau_d s}. \tag{12}
\end{equation}
This is a stationary curve identity at actual speed $s$. It does not prove
that native transient waves propagate at this $w$, that advice is instantly
realized, or that a separately served lane stream has speed $U(g_j)$.
SUMO's documented headway variations and insertion effects matter to
any later admission. No runtime capacity was measured in this review.

\newpage
\section*{4. Recovery scale and joint-domain implications}
For $u<u_R=U(R)$ the isolated active-profile duration and integrated loss are
\begin{equation}
 \Theta(u)=\frac{(w+u_R)^2-(w+u)^2}{2wa_r},\qquad
 W(u)=\int_0^{\Theta(u)}D(\theta;u)\,d\theta
      =\frac{K}{w+u_R}\frac{(u_R-u)^2}{2a_r}. \tag{13}
\end{equation}
For $u\geq u_R$, both are zero. While the profile is active,
\begin{equation}
 \frac{\partial D}{\partial\theta}=-\frac{Kw^2a_r}{S^3},\quad
 \frac{\partial D}{\partial u}=-\frac{Kw(w+u)}{S^3},\quad
 \frac{\partial D}{\partial a_r}=-\frac{Kw^2\theta}{S^3}. \tag{14}
\end{equation}
At fixed admission times and marks, stronger acceleration reduces each
surviving profile and the account in (3). In the coupled traffic model,
admissions and marks also change. This is not a theorem that complete TST
or optimal advice varies monotonically with acceleration.

$W$ has units veh-eq and $\Theta$ has units s. They are profile descriptors,
not new current E charges. They constrain different aspects of recovery:
a total loss alone confounds low insertion speed with slow acceleration;
independent marks and time-resolved tails are more informative. O-R's
retirement may truncate an older profile before acceleration removes it,
so the isolated area is not the integrated account of an arbitrary history.

\textbf{Joint admission.} Require one common parameter set that explains
preparation, actual speed marks, mixed discharge and recovery, and lies in
the existing full-advice forward domain. In particular retain
\begin{equation}
 v_+r<R\leq C(v_-),\qquad M<k_jb,\qquad
 c_0-\beta(1+v_+/w)>0,\qquad a_r>0. \tag{15}
\end{equation}
Other H1/T1 and far-lane assumptions remain as previously declared.
$\chi=\beta/[C_\star v_\star]$ is anchored once, never refitted by advice
or search box. $m$ remains the selected physical normalization. (15) is
mathematical admission, not an observed physical range. A good objective
fit cannot replace the mechanism checks. No numerical value, tolerance,
positive policy result or theory-only primary is selected here.

\textbf{Queue-scope clarification.} The earlier E non-requeue conclusion
uses the monotone H1 pulse, free-speed/recovery rules and retained domain.
Finiteness alone is insufficient: a bursty later E wave can requeue in a
different model or a native finite cohort. Such an event cannot by itself
be used to reject the independent reference.

\textbf{Source ownership.} Carey Section 5 motivates receiving work;
Leclercq Sections 2 and 3.2 constrain the recovery and speed interpretations;
official SUMO road-capacity documentation supports only the stationary
correspondence. Equations (4)--(10), the admission order and joined-law
qualification are our deductions. No source validates the complete node.
See the accompanying viewed-reference list and physical admission contract.
Applied MP6 and Step 3 remain open; Step 4 is unopened.
\end{document}
